\documentclass{article}
\usepackage{amsmath,amssymb}
\usepackage{algorithm}
\usepackage{algorithmic}

\begin{document}

\section{Policy Gradient Theorem}

Steps to derive the policy gradient theorem:

\begin{enumerate}
    \item Define a policy \( \pi_\theta(a|s) \) parameterized by \( \theta \), which gives the probability of taking action \( a \) in state \( s \).
    \item Define a a state-action trajectory \( \tau = (s_0, a_0, s_1, a_1, \ldots s_T, a_T) \) generated by following the policy \( \pi_\theta \).
    \item Define the return of a trajectory \( G (\tau) \) as the discounted sum of rewards observed along the trajectory:
    \begin{equation}
    G(\tau) = \sum_{t} \gamma^t r (s_t, a_t)
    \end{equation}
    where \( \gamma \) is the discount factor.
    \item Start with the definition of the objective function \( J(\theta) \) that we want to maximize, which is the expected return when following the policy parameterized by \( \theta \):
    \begin{equation}
    J(\theta) = \mathbb{E}_{\tau \sim \pi_\theta} [G (\tau)] = \int_{\tau} P(\tau ; \theta) G(\tau) d\tau
    \end{equation}
\end{enumerate}

\section{REINFORCE Algorithm}


\section{The variance problem of REINFORCE}
\quad The REINFORCE algorithm is a Monte Carlo policy gradient method that estimates the policy gradient using the returns observed in complete episodes.
While REINFORCE is a simple and intuitive algorithm, it suffers from high variance in the policy gradient estimates, which can lead to slow convergence and unstable learning.
The high variance arises because the returns \( G(\tau) \) can vary widely across different trajectories, especially in environments with long episodes or sparse rewards. 
This means that the policy updates can be based on very noisy estimates of the expected return, which can cause the learning process to be inefficient and unstable.


\section{Actor-Critic Methods}


\end{document}